\chapter[Vierbein emergence]{Condensation, Vierbein Emergence, and the Seiberg--Witten Map}
\label{ch:vierbein-emergence}

\paragraph{Learning outcomes.}
After this chapter the reader should be able to:
\begin{enumerate}
  \item explain why a strictly constant vacuum
        \(\langle\Phi^I{}_A\rangle=v\delta^I{}_A\) cannot, by itself,
        produce a non-degenerate vierbein from
        \(\langle\partial_\mu\Phi\rangle\);
  \item write the polar decomposition
        \(\Phi=(v+\sigma)\Lambda\), state the polar-derivative
        prescription of the source notes as an unverified
        proposal that does not supply a frame (Problem~C), and
        state the frame of Repair~II;
  \item reconstruct \(g_{\mu\nu}=\eta_{AB}e^A_\mu e^B_\nu\) and state
        the non-degeneracy condition;
  \item use the leading Seiberg--Witten map to relate
        \(\Theta^{\mu\nu}\) to a correction of the emergent metric;
  \item enter Chapter~5 with a definite background about which to
        expand \(\delta\Phi\).
\end{enumerate}

Chapter~3 introduced the field and a preferred vacuum. This chapter
turns that vacuum into geometry conditionally: it posits the Repair~II
spacetime-dependent condensate as the frame, since the Repair~I
prescription does not supply one. No Einstein equation is derived
here. The only output is a composite metric, plus the dictionary that
connects it to the non-commutative parameter \(\Theta^{\mu\nu}\).

\section{The constant-vacuum obstruction}
\label{sec:ch4-obstruction}

The representative vacuum of Assumption~\ref{ass:frame-lock} is
\begin{equation}
  \label{eq:const-vev}
  \langle\Phi^I{}_A\rangle=v\,\delta^I{}_A.
\end{equation}
If this expectation value is independent of \(x\), then
\begin{equation}
  \label{eq:vanishing-frame}
  \langle\partial_\mu\Phi^{IA}\rangle=0.
\end{equation}
The definition advertised in the source notes,
\begin{equation}
  \label{eq:naive-vierbein}
  \mathcal{E}^A{}_{\mu I}(x)
  \equiv
  \langle\partial_\mu\Phi^{IA}(x)\rangle,
\end{equation}
then yields the zero tensor. A zero tensor is not a frame.

This is not a notational quibble. A textbook that passes from
\eqref{eq:const-vev} to a non-degenerate \(g_{\mu\nu}\) without an
intermediate step is asserting an identity that is false. There are
two candidate repairs. Only Repair~II is carried through in this book;
they must be kept distinct.

\paragraph{Repair I: the proposed polar-derivative prescription.}
Write the field in polar form, as in the ultraviolet note,
\begin{equation}
  \label{eq:polar}
  \Phi^{IA}(x)=\bigl(v+\sigma(x)\bigr)\,\Lambda^A{}_I(x),
\end{equation}
where \(\sigma\) is a radial singlet and \(\Lambda(x)\) takes values
in the broken coset that rotates internal indices into tangent
indices. Even if \(v\) is constant, \(\partial_\mu\Phi\) need not
vanish:
\begin{equation}
  \label{eq:polar-deriv}
  \partial_\mu\Phi^{IA}
  =(\partial_\mu\sigma)\,\Lambda^A{}_I
   +(v+\sigma)\,\partial_\mu\Lambda^A{}_I.
\end{equation}
A vacuum with \(\langle\sigma\rangle=0\) and a slowly varying
\(\langle\Lambda\rangle\) has \(\langle\partial_\mu\Lambda\rangle\ne0\),
but this is a candidate frame only in the sense of the source
notes; see the discussion after Repair~II.

\paragraph{Repair II: identify the vacuum itself with a frame.}
Promote the representative vacuum to a spacetime-dependent but
still ``pure Higgs'' configuration
\begin{equation}
  \label{eq:higgs-frame}
  \langle\Phi^I{}_A(x)\rangle=v\,e^I{}_A(x),
\end{equation}
and treat \(e^I{}_A\) as the fundamental condensate profile. In this
repair, equation~\eqref{eq:naive-vierbein} is abandoned as a
definition. The vierbein \emph{is} the normalized vacuum,
\begin{equation}
  \label{eq:vev-is-frame}
  e^A_\mu(x)
  \equiv
  \frac1v\,\langle\Phi^{A}{}_\mu(x)\rangle,
\end{equation}
after the internal and tangent bases have been identified.

Repair~II is simpler and is the standard Higgs-gravity pattern.
Repair~I is what the notes write when they insist on
\eqref{eq:naive-vierbein}. The two are not equivalent. Repair~II
defines a frame directly from the condensate and, wherever
Assumption~\ref{ass:nondeg} holds, supplies a non-degenerate
\(e^A{}_\mu\). Repair~I, as written in
Definition~\ref{def:vierbein}, leaves the index \(I\) uncontracted,
and Appendix~\ref{app:status} (Section~\ref{sec:status-C}) shows that no
Lorentz-equivariant linear contraction of \(\partial_\mu\Lambda\)
produces a vierbein fluctuation. Definition~\ref{def:vierbein} is
therefore kept only as a record of the prescription in the source
notes; it is not a constructed frame. In the rest of the book the
frame is the Repair~II frame~\eqref{eq:vev-is-frame}, or a
non-degenerate coframe assumed separately
(Assumption~\ref{ass:nondeg}); statements that use the explicit form
\eqref{eq:official-vierbein} are conditional on a repair of Repair~I
that this book does not supply.

\begin{definition}[Repair~I prescription of the source notes]
  \label{def:vierbein}
  In a polar configuration~\eqref{eq:polar} with
  \(\langle\sigma\rangle=0\), define the three-index object
  \begin{equation}
    \label{eq:official-vierbein}
    \mathcal{E}^{A}{}_{\mu I}(x)
    \equiv
    \bigl\langle\partial_\mu\Lambda^{A}{}_I(x)\bigr\rangle .
  \end{equation}
  By~\eqref{eq:polar-deriv} this is
  \(v^{-1}\langle\partial_\mu\Phi^{IA}\rangle\) at \(\sigma=0\)
  (with \(\langle\partial_\mu\sigma\rangle=0\)), which is the
  prescription the source notes write as a vierbein, with the
  internal index \(I\) left uncontracted and no identification of
  \(I\) with a tangent index used. It is not a vierbein: it carries
  the extra index \(I\), and no Lorentz-equivariant linear contraction
  of it yields a vierbein fluctuation (Appendix~\ref{app:status},
  Section~\ref{sec:status-C}). The factor \(1/v\) is the normalisation
  the notes use when \(\Phi\) has the engineering dimension of \(v\);
  if a formula in the notes omits \(1/v\), it is using units \(v=1\).
\end{definition}

\begin{remark}[The two repairs do not agree]
The expansion
\(\Lambda^A{}_I=\delta^A{}_I+v^{-1}\omega^A{}_{I\mu}x^\mu+\cdots\)
makes \(\partial_\mu\Lambda^A{}_I\) a constant three-index tensor,
antisymmetric in \((A,I)\). It lies in
\(V\otimes\mathfrak{so}(3,1)\), while a vierbein perturbation lies
in \(V\otimes V\); the two share no irreducible summand
(Appendix~\ref{app:status}, Section~\ref{sec:status-C}). Repair~I
therefore does not reproduce Repair~II at first order. A
corrected Repair~I would need additional background tensors or an
independent coframe, which is part of Problem~C.
\end{remark}

\section{The composite metric}
\label{sec:ch4-metric}

\begin{definition}[Composite metric]
  \label{def:composite-metric}
  Given a non-degenerate frame \(e^A_\mu\), taken from
  Repair~II,~\eqref{eq:vev-is-frame}, or supplied separately as an
  assumption (Assumption~\ref{ass:nondeg}),
  \begin{equation}
    \label{eq:g-from-e}
    g_{\mu\nu}(x)
    \equiv
    \eta_{AB}\,e^A_\mu(x)\,e^B_\nu(x).
  \end{equation}
\end{definition}

Non-degeneracy is not automatic.

\begin{assumption}[Non-degenerate condensate frame]
  \label{ass:nondeg}
  On the region of spacetime under consideration, \(e^A_\mu\) has
  non-vanishing determinant. Equivalently, the four covectors
  \(e^A{}_\mu\mathrm{d}x^\mu\) are linearly independent.
\end{assumption}

If Assumption~\ref{ass:nondeg} fails, \(g_{\mu\nu}\) is degenerate and
there is no Lorentzian geometry to discuss. Vanishing of
\(\det e\) is a physical possibility---a condensate defect, a
region where the frame-locking vacuum has not formed---not a gauge
artifact. This book restricts attention to the condensed phase in
which Assumption~\ref{ass:nondeg} holds.

The inverse metric and inverse frame are defined in the usual way
once non-degeneracy is granted,
\begin{equation}
  \label{eq:inverses}
  g^{\mu\nu}e^A_\mu e^B_\nu=\eta^{AB},\qquad
  e^\mu_A e^A_\nu=\delta^\mu_\nu.
\end{equation}
Local Lorentz transformations act on the frame index,
\begin{equation}
  \label{eq:llts}
  e^A_\mu\;\longmapsto\;\Lambda^A{}_B e^B_\mu,
\end{equation}
and leave \(g_{\mu\nu}\) invariant. Diffeomorphisms act on the
coordinate index. These two redundancies are the reason a composite
metric can have the same geometric content as a primitive metric,
provided the Goldstone counting of Chapter~10 works.

\section{What kind of object \(e^A_\mu\) is}
\label{sec:ch4-status}

At this stage \(e^A_\mu\) is a condensate profile, not a solution of
Einstein's equation. It is closer to a magnetization than to a
gravitational field. Three consequences follow.

\begin{enumerate}
  \item There is no a priori reason for \(e^A_\mu\) to satisfy the
        torsion-free condition~\eqref{eq:spin-connection}. If torsion
        is to vanish, that is a dynamical statement about the
        connection that eats the Goldstones, not a definition.
  \item There is no a priori reason for small disturbances of
        \(g_{\mu\nu}\) to propagate as massless spin-2 waves. That is
        Chapter~7.
  \item The auxiliary scaffold \(g^{(0)}_{\mu\nu}\) of
        Presentation~B is still present in the microscopic action.
        It must not be confused with~\eqref{eq:g-from-e}. Heat-kernel
        calculations in Chapter~8 are performed on a background that
        is eventually identified with the emergent metric; making
        that identification precise is part of the coefficient map.
\end{enumerate}

\section{Seiberg--Witten map and the \(\Theta\) correction}
\label{sec:ch4-SW}

Presentation~A still has a non-commutative parameter
\(\Theta^{\mu\nu}\). After a commutative manifold has been introduced,
the open-string / non-commutative-gauge-theory dictionary of Seiberg
and Witten relates an ordinary metric \(g_{\mu\nu}\) and a two-form
background to an open-string effective metric. The source notes write
the leading relation
\begin{equation}
  \label{eq:SW}
  G_{\mu\nu}(x)
  =\bigl((g+\Theta)^{-1}\bigr)_{\mu\nu}
   +\mathcal{G}_{\mu\nu}\bigl(\langle\hat X\rangle\bigr).
\end{equation}

Two clarifications are required before~\eqref{eq:SW} is used.

\paragraph{Matrix versus form.}
As written, \(g+\Theta\) adds a symmetric tensor to an antisymmetric
tensor. The honest object is the mixed tensor
\begin{equation}
  \label{eq:g-plus-Theta}
  (g+\Theta)_{\mu}{}^{\nu}
  =g_{\mu\rho}\eta^{\rho\nu}+\Theta_{\mu}{}^{\nu},
\end{equation}
or, equivalently, the linear map \(g+\Theta:T^*M\to TM\) after one
index of \(\Theta\) is raised with the auxiliary metric. The inverse
in~\eqref{eq:SW} is the inverse of that map. In components, for
constant \(\Theta\) on a flat auxiliary background,
\begin{equation}
  \label{eq:SW-flat}
  \bigl((g+\Theta)^{-1}\bigr)_{\mu\nu}
  =g_{\mu\nu}
   -g_{\mu\alpha}\Theta^{\alpha\beta}g_{\beta\nu}
   +O(\Theta^2).
\end{equation}
The first correction is symmetric if one takes the symmetric part,
which is the part that can enter a metric. The antisymmetric remainder
is a two-form background, not a correction to \(g_{\mu\nu}\).

\paragraph{The extra term \(\mathcal{G}_{\mu\nu}\).}
The piece \(\mathcal{G}_{\mu\nu}(\langle\hat X\rangle)\) in
\eqref{eq:SW} is the backreaction of the matrix condensate on the
emergent geometry, beyond the kinematic \(\Theta\) correction. At
this stage of the book it is a name, not a formula. Chapters~8
and~11 compute pieces of it: the Einstein--Hilbert term from the
heat kernel, and the entropic stress tensor from mutual information.
Until those calculations exist, \(\mathcal{G}_{\mu\nu}\) must not be
assigned a phenomenology.

\begin{definition}[Working infrared metric]
  \label{def:working-metric}
  Unless a higher-order \(\Theta\) correction is being tracked
  explicitly, ``the emergent metric'' means~\eqref{eq:g-from-e}.
  Equation~\eqref{eq:SW} is the next-order kinematic correction
  plus a placeholder for dynamical backreaction.
\end{definition}

\section{A linearised frame around Minkowski spacetime}
\label{sec:ch4-minkowski}

The fluctuation chapters need a background. Choose coordinates and a
local Lorentz gauge so that the condensed frame is the trivial frame
plus a small disturbance,
\begin{equation}
  \label{eq:linear-frame}
  e^A_\mu=\delta^A_\mu+\tfrac12 h^A{}_\mu+\omega^A{}_\mu,
\end{equation}
where \(h_{\mu\nu}\equiv\eta_{AB}\delta^A_\mu h^B{}_\nu\) is symmetric
and \(\omega_{\mu\nu}\) is antisymmetric. Then
\begin{equation}
  \label{eq:linear-metric}
  g_{\mu\nu}
  =\eta_{\mu\nu}+h_{\mu\nu}+O(h^2,\omega^2,h\omega).
\end{equation}
The antisymmetric piece \(\omega\) is a local Lorentz Goldstone. If
Chapter~10 operates as advertised, \(\omega\) is eaten and does not
propagate. The symmetric piece \(h_{\mu\nu}\) is the candidate
graviton, still un-decomposed and possibly massive. Chapter~5 splits
it; Chapter~7 asks whether its mass vanishes.

This expansion assumes that a Minkowski condensate exists: a vacuum
in which \(e^A_\mu=\delta^A_\mu\) globally, or in a region large
compared with the scale of later experiments. Existence of that
vacuum is part of Assumption~\ref{ass:nondeg}, not a new principle.

\section{Signature and the role of \(\eta_{AB}\)}
\label{sec:ch4-signature}

Equation~\eqref{eq:g-from-e} borrows its signature from \(\eta_{AB}\).
Nothing in the potential~\eqref{eq:V} selected
\((-,+,+,+)\) rather than \((+,+,+,+)\). A Euclidean condensate is a
priori as natural as a Lorentzian one.

\begin{assumption}[Lorentzian condensate]
  \label{ass:lorentzian}
  The locked vacuum is Lorentzian:
  \(\eta_{AB}=\mathrm{diag}(-1,+1,+1,+1)\). Dynamical emergence of
  this signature from a Euclidean matrix model is not proved in the
  source notes and is not pretended here. Chapter~10 shows how local
  \(\mathrm{SO}(3,1)\) can arise \emph{given} this \(\eta_{AB}\).
\end{assumption}

A reader looking for a derivation of Lorentzian signature from
nothing will not find it in this book. What the book does derive,
contingent on Assumption~\ref{ass:lorentzian}, is the dynamics of
the resulting Lorentzian metric.

\section{Worked example: a weak, static frame perturbation}
\label{sec:ch4-example}

Take the Repair~II frame, or any assumed non-degenerate frame, with
the constant background \(e^A{}_\mu=\delta^A{}_\mu\), and a small
coordinate shift \(x^\mu\to x^\mu+\zeta^\mu(x)\). The frame changes by
its Lie derivative,
\begin{equation}
  \label{eq:example-e}
  \delta e^A{}_\mu
  =\zeta^\nu\partial_\nu\delta^A{}_\mu
   +\partial_\mu\zeta^\nu\,\delta^A{}_\nu
  =\partial_\mu\zeta^A .
\end{equation}
The associated metric perturbation is
\begin{equation}
  \label{eq:example-h}
  h_{\mu\nu}
  =\partial_\mu\zeta_\nu+\partial_\nu\zeta_\mu
\end{equation}
to first order. This configuration is geometry in name only: it is a
coordinate artifact. A physical graviton must be a disturbance that
cannot be written as~\eqref{eq:example-h}. Isolating that disturbance
is the point of the irreducible decomposition in Chapter~5.

The example demonstrates diffeomorphism redundancy of a Repair~II
(or assumed) frame. It does not derive a frame from \(\Phi\) by
Definition~\ref{def:vierbein}.

\section{What this chapter has not done}
\label{sec:ch4-not-done}

The following statements are \emph{not} theorems of Chapter~4.

\begin{itemize}
  \item \(g_{\mu\nu}\) obeys Einstein's equation.
  \item \(h_{\mu\nu}\) contains two massless TT polarizations.
  \item \(\Geff\) has the value~\eqref{eq:Geff-again}.
  \item Torsion vanishes.
  \item The signature is dynamical.
  \item \(\mathcal{G}_{\mu\nu}\) has been computed.
\end{itemize}

What has been done is narrower and indispensable: a composite metric
has been defined from a non-degenerate condensate frame
(Assumption~\ref{ass:nondeg}); the constant-vacuum obstruction
is addressed only by positing, in Repair~II, a spacetime-dependent
condensate that is identified with the frame~\eqref{eq:vev-is-frame}
(no such profile is derived), while the Repair~I formula of
Definition~\ref{def:vierbein} is not a constructed frame
(Problem~C, Appendix~\ref{app:status}); and the linearized
Minkowski background used in every later proof has been fixed.

\section{Exercises}
\label{sec:ch4-exercises}

\begin{enumerate}
  \item \textbf{Zero frame.}
    Prove that a strictly constant \(\langle\Phi\rangle\) yields
    \(\langle\partial_\mu\Phi\rangle=0\). This is the obstruction of
    Section~\ref{sec:ch4-obstruction} as a two-line exercise.
  \item \textbf{Polar derivative.}
    Derive~\eqref{eq:polar-deriv} from~\eqref{eq:polar}. Under what
    condition on \(\Lambda\) is \(e^A_\mu\) proportional to a
    Maurer--Cartan form \(\Lambda^{-1}\partial_\mu\Lambda\)?
  \item \textbf{Non-degeneracy.}
    Let \(e^A_\mu=\mathrm{diag}(1,1,1,\varepsilon)\). Compute
    \(g_{\mu\nu}\) and \(\det g\). What fails as
    \(\varepsilon\to 0\)? Interpret the limit as a condensate defect.
  \item \textbf{Symmetric part of the \(\Theta\) correction.}
    From~\eqref{eq:SW-flat}, extract the symmetric \(O(\Theta)\)
    correction to \(g_{\mu\nu}\) and the antisymmetric remainder.
    Which object can enter geodesic motion at this order?
  \item \textbf{Linearized metric.}
    Starting from~\eqref{eq:linear-frame}, derive
    \eqref{eq:linear-metric} and show that \(\omega_{AB}\) drops out
    of \(g_{\mu\nu}\) at first order. Where does \(\omega\) re-enter
    at second order?
  \item \textbf{Pure gauge.}
    Show that~\eqref{eq:example-h} is the linearized diffeomorphism
    action on \(\eta_{\mu\nu}\). Write one explicit \(h_{\mu\nu}\)
    that cannot be written in this form. (Any TT wave will do.)
  \item \textbf{Auxiliary versus emergent.}
    In Presentation~B the microscopic action uses \(g^{(0)}\). In
    this chapter the geometry is \(g\). Write one sentence stating
    when it is legitimate to replace \(g^{(0)}\) by \(g\) inside a
    heat kernel, and what error term that replacement drops. The
    full answer is Chapter~8; the sentence is due now.
\end{enumerate}

\section{Bridge to Chapter~5}
\label{sec:ch4-bridge}

The background is a condensed frame, taken Minkowski for the
linearized theory,
\begin{equation}
  \label{eq:bg-next}
  \bar e^A_\mu=\delta^A_\mu,
  \qquad
  \bar g_{\mu\nu}=\eta_{\mu\nu}.
\end{equation}
The next object is the fluctuation \(\delta\Phi^I{}_A\) about the
polar vacuum~\eqref{eq:polar}. Chapter~5 decomposes that fluctuation
under local Lorentz transformations into a conformal piece, a vector
piece, and a symmetric trace-free tensor \(T_{\mu\nu}\). Only after
that decomposition is it meaningful to ask which piece might be a
graviton.